%% 自動生成（tools/mkchecklist.py）。各章の「ふりかえり」の項目を集めたもの
\begingroup\footnotesize
\def\bkcb{\tikz[baseline=-0.5ex]\draw[\bkTF, line width=0.7pt, rounded corners=0.5pt] (0,-0.3em) rectangle (0.8em,0.5em);}
\begin{longtblr}{colspec={Q[c,m,wd=9mm] X[l,m] Q[c,m,wd=10mm] Q[c,m,wd=9mm] Q[c,m,wd=9mm]}, width=\linewidth, rowhead=1, colsep=3pt,
  row{1}={bg=bkVol!14, font=\sffamily\bfseries\scriptsize}, hlines={0.4pt,bkLine}, rowsep=2.3pt}
\R{章}{しょう} & できるようになったこと & ページ & 1\R{回}{かい}め & 2\R{回}{かい}め \\
{\sffamily\bfseries\color{\bkTT}3-1} & \R{式}{しき}の\R{中}{なか}の\R{変数}{へん|すう}と\R{定数}{てい|すう}を\R{見分}{み|わ}けられる。 & \pageref{bk:3-1:hensu} & \bkcb & \bkcb \\
 & \R{物理定数}{ぶつ|り|てい|すう}と、\R{問題}{もん|だい}ごとに\R{決}{き}まる\R{数}{かず}のちがいを\R{説明}{せつ|めい}できる。 & \pageref{bk:3-1:butsuri} & \bkcb & \bkcb \\
 & \R{変域}{へん|いき}を\R{不等号}{ふ|とう|ごう}で\R{書}{か}ける。 & \pageref{bk:3-1:heniki} & \bkcb & \bkcb \\
{\sffamily\bfseries\color{\bkTT}3-2} & \R{原点}{げん|てん}と\R{軸}{じく}を\R{決}{き}めて、\R{点}{てん}をかいたり\R{座標}{ざ|ひょう}を\R{読}{よ}んだりできる。 & \pageref{bk:3-2:zahyo} & \bkcb & \bkcb \\
 & マイナスの\R{座標}{ざ|ひょう}、4つの\R{区画}{く|かく}、\R{点}{てん}を\R{動}{うご}かした\R{座標}{ざ|ひょう}がわかる。 & \pageref{bk:3-2:ugokasu} & \bkcb & \bkcb \\
 & \R{絶対値}{ぜっ|たい|ち}を\R{使}{つか}って、2\R{点}{てん}の\R{距離}{きょ|り}を\R{求}{もと}められる。 & \pageref{bk:3-2:kyori} & \bkcb & \bkcb \\
{\sffamily\bfseries\color{\bkTT}3-3} & \R{関数}{かん|すう}を\R{見分}{み|わ}け、\R{表}{ひょう}・グラフ・\R{式}{しき}で\R{表}{あらわ}せる。 & \pageref{bk:3-3:kansu} & \bkcb & \bkcb \\
 & $f(x)$ を\R{読}{よ}み、$f(3)$ や $f(a+h)$ を\R{計算}{けい|さん}できる。 & \pageref{bk:3-3:fx} & \bkcb & \bkcb \\
 & 2つの\R{関数}{かん|すう}をつないだ\R{合成}{ごう|せい}\R{関数}{かん|すう}を\R{作}{つく}れる。 & \pageref{bk:3-3:gousei} & \bkcb & \bkcb \\
{\sffamily\bfseries\color{\bkTT}3-4} & \R{軸}{じく}と\R{目}{め}もりを\R{決}{き}めて、はかった\R{点}{てん}からグラフをかける。 & \pageref{bk:3-4:jiku} & \bkcb & \bkcb \\
 & $x$-$t$ グラフと $v$-$t$ グラフを\R{読}{よ}み\R{分}{わ}けられる。 & \pageref{bk:3-4:xt} & \bkcb & \bkcb \\
 & \R{表}{ひょう}にない\R{値}{あたい}を\R{線}{せん}から\R{読}{よ}み、\R{外挿}{がい|そう}の\R{注意}{ちゅう|い}を\R{説明}{せつ|めい}できる。 & \pageref{bk:3-4:naisou} & \bkcb & \bkcb \\
{\sffamily\bfseries\color{\bkTT}3-5} & \R{傾}{かたむ}きを「たての\R{変化量}{へん|か|りょう}÷\R{横}{よこ}の\R{変化量}{へん|か|りょう}」で\R{求}{もと}め、\R{正}{せい}・\R{負}{ふ}・0がわかる。 & \pageref{bk:3-5:katamuki} & \bkcb & \bkcb \\
 & $y=ax+b$ の\R{傾}{かたむ}きと\R{切片}{せっ|ぺん}を\R{読}{よ}み、$x$-$t$ グラフの\R{傾}{かたむ}きを\R{速}{はや}さとして\R{使}{つか}える。 & \pageref{bk:3-5:ichiji} & \bkcb & \bkcb \\
 & 1\R{点}{てん}と\R{傾}{かたむ}きから、$y-y_1=m(x-x_1)$ で\R{直線}{ちょく|せん}の\R{式}{しき}を\R{作}{つく}れる。 & \pageref{bk:3-5:tenkatamuki} & \bkcb & \bkcb \\
{\sffamily\bfseries\color{\bkTT}3-6} & \R{交点}{こう|てん}が、2つの\R{式}{しき}を\R{同時}{どう|じ}にみたす\R{組}{くみ}だとわかる。 & \pageref{bk:3-6:kouten} & \bkcb & \bkcb \\
 & \R{代入法}{だい|にゅう|ほう}と\R{加減法}{か|げん|ほう}で\R{解}{と}き、\R{両方}{りょう|ほう}の\R{式}{しき}で\R{確}{たし}かめられる。 & \pageref{bk:3-6:dainyuu} & \bkcb & \bkcb \\
 & \R{解}{かい}がない・\R{無数}{む|すう}のわけを\R{説明}{せつ|めい}し、\R{出会}{で|あ}い・\R{買}{か}い\R{物}{もの}・\R{回路}{かい|ろ}の\R{問題}{もん|だい}を\R{解}{と}ける。 & \pageref{bk:3-6:heikou} & \bkcb & \bkcb \\
{\sffamily\bfseries\color{\bkTT}3-7} & \R{双曲線}{そう|きょく|せん}をかき、\R{軸}{じく}にくっつかないわけを\R{言}{い}える。 & \pageref{bk:3-7:jiku} & \bkcb & \bkcb \\
 & $a<0$ のグラフを\R{読}{よ}み、\R{変化}{へん|か}の\R{割合}{わり|あい}を\R{計算}{けい|さん}できる。 & \pageref{bk:3-7:wariai} & \bkcb & \bkcb \\
 & $xy$ が\R{一定}{いっ|てい}を\R{見}{み}つけ、\R{式}{しき}を\R{作}{つく}って\R{使}{つか}える。 & \pageref{bk:3-7:pv} & \bkcb & \bkcb \\
{\sffamily\bfseries\color{\bkTT}3-8} & かっこのある\R{式}{しき}を\R{展開}{てん|かい}し、\R{面積}{めん|せき}の\R{図}{ず}で\R{説明}{せつ|めい}できる。 & \pageref{bk:3-8:kumiawase} & \bkcb & \bkcb \\
 & \R{恒等式}{こう|とう|しき}の\R{係数}{けい|すう}を、2\R{通}{とお}りの\R{方法}{ほう|ほう}で\R{決}{き}められる。 & \pageref{bk:3-8:koutou} & \bkcb & \bkcb \\
 & \R{平方完成}{へい|ほう|かん|せい}して、いちばん\R{小}{ちい}さい・\R{大}{おお}きい\R{値}{あたい}を\R{読}{よ}める。 & \pageref{bk:3-8:heihou} & \bkcb & \bkcb \\
{\sffamily\bfseries\color{\bkTT}3-9} & 2\R{乗}{じょう}に\R{比例}{ひ|れい}を\R{見}{み}つけ、$y=ax^2$ の\R{式}{しき}とグラフにできる。 & \pageref{bk:3-9:nijo} & \bkcb & \bkcb \\
 & $a$ とグラフの\R{形}{かたち}の\R{関係}{かん|けい}や、\R{変化}{へん|か}の\R{割合}{わり|あい}を\R{説明}{せつ|めい}できる。 & \pageref{bk:3-9:a} & \bkcb & \bkcb \\
 & \R{平方完成}{へい|ほう|かん|せい}で\R{頂点}{ちょう|てん}を、\R{範囲}{はん|い}から\R{最大}{さい|だい}・\R{最小}{さい|しょう}を\R{求}{もと}められる。 & \pageref{bk:3-9:heihou} & \bkcb & \bkcb \\
{\sffamily\bfseries\color{\bkTT}3-10} & $x^2=k$ の\R{形}{かたち}や\R{因数分解}{いん|すう|ぶん|かい}で、2\R{次方程式}{じ|ほう|てい|しき}を\R{解}{と}ける。 & \pageref{bk:3-10:insu} & \bkcb & \bkcb \\
 & \R{解}{かい}の\R{公式}{こう|しき}を\R{作}{つく}り、$\sqrt{\phantom{2}}$ の\R{計算}{けい|さん}をして\R{解}{と}ける。 & \pageref{bk:3-10:kousiki} & \bkcb & \bkcb \\
 & \R{判別式}{はん|べつ|しき}で\R{解}{かい}の\R{個数}{こ|すう}を\R{調}{しら}べ、\R{場面}{ば|めん}に\R{合}{あ}う\R{解}{かい}を\R{選}{えら}べる。 & \pageref{bk:3-10:hanbetsu} & \bkcb & \bkcb \\
{\sffamily\bfseries\color{\bkTT}3-11} & 「ふえる\R{量}{りょう}が、いまある\R{量}{りょう}に\R{比例}{ひ|れい}する」\R{変化}{へん|か}を、\R{表}{ひょう}・\R{式}{しき}・グラフにできる。 & \pageref{bk:3-11:hirei} & \bkcb & \bkcb \\
 & \R{指数}{し|すう}のきまり $a^{x+h}=a^x\times a^h$、$(a^m)^n=a^{mn}$、$a^0=1$ を\R{使}{つか}える。 & \pageref{bk:3-11:kimari} & \bkcb & \bkcb \\
 & $x^{\frac{1}{2}}=\sqrt{x}$ のわけを\R{説明}{せつ|めい}できる。 & \pageref{bk:3-11:bunsu} & \bkcb & \bkcb \\
{\sffamily\bfseries\color{\bkTT}3-12} & 10を\R{何乗}{なん|じょう}するかでけたを\R{数}{かぞ}え、$\log_ab$ の\R{意味}{い|み}を\R{説明}{せつ|めい}できる。 & \pageref{bk:3-12:log} & \bkcb & \bkcb \\
 & \R{対数}{たい|すう}のきまりを、\R{指数}{し|すう}のきまりから\R{作}{つく}って\R{使}{つか}える。 & \pageref{bk:3-12:kimari} & \bkcb & \bkcb \\
 & \R{片対数}{かた|たい|すう}グラフやデシベルを\R{使}{つか}える。 & \pageref{bk:3-12:kata} & \bkcb & \bkcb \\
{\sffamily\bfseries\color{\bkTT}3-13} & \R{何回}{なん|かい}もはかった\R{値}{あたい}の\R{平均}{へい|きん}を、$\bar{x}$ を\R{使}{つか}って\R{求}{もと}められる。 & \pageref{bk:3-13:heikin} & \bkcb & \bkcb \\
 & \R{度数分布表}{ど|すう|ぶん|ぷ|ひょう}とヒストグラムを\R{作}{つく}り、\R{読}{よ}みとれる。 & \pageref{bk:3-13:dosuu} & \bkcb & \bkcb \\
 & \R{中央値}{ちゅう|おう|ち}・\R{最頻値}{さい|ひん|ち}・\R{四分位数}{し|ぶん|い|すう}を\R{求}{もと}め、\R{箱}{はこ}ひげ\R{図}{ず}をかける。 & \pageref{bk:3-13:daihyou} & \bkcb & \bkcb \\
{\sffamily\bfseries\color{\bkTT}3-14} & \R{偏差}{へん|さ}を\R{求}{もと}め、\R{合計}{ごう|けい}が0になるわけを\R{説明}{せつ|めい}できる。 & \pageref{bk:3-14:hensa} & \bkcb & \bkcb \\
 & \R{分散}{ぶん|さん}と $\sigma$ を\R{求}{もと}め、ばらつきをくらべられる。 & \pageref{bk:3-14:sigma} & \bkcb & \bkcb \\
 & \R{平均}{へい|きん}$\pm\sigma$ と $\tfrac{\sigma}{\sqrt{n}}$ の\R{意味}{い|み}を\R{説明}{せつ|めい}できる。 & \pageref{bk:3-14:nanawari} & \bkcb & \bkcb \\
{\sffamily\bfseries\color{\bkTT}3-15} & \R{散布図}{さん|ぷ|ず}をかき、\R{相関}{そう|かん}の\R{向}{む}きと\R{強}{つよ}さを\R{見分}{み|わ}けられる。 & \pageref{bk:3-15:sanpu} & \bkcb & \bkcb \\
 & \R{共分散}{きょう|ぶん|さん}を\R{計算}{けい|さん}し、その\R{符号}{ふ|ごう}の\R{意味}{い|み}を\R{説明}{せつ|めい}できる。 & \pageref{bk:3-15:kyobunsan} & \bkcb & \bkcb \\
 & \R{相関係数}{そう|かん|けい|すう} $r$ を\R{求}{もと}め、\R{相関}{そう|かん}と\R{因果}{いん|が}を\R{区別}{く|べつ}できる。 & \pageref{bk:3-15:r} & \bkcb & \bkcb \\
{\sffamily\bfseries\color{\bkTT}3-16} & ずれの2\R{乗}{じょう}の\R{和}{わ}で、どの\R{直線}{ちょく|せん}がよく\R{合}{あ}うかくらべられる。 & \pageref{bk:3-16:zure} & \bkcb & \bkcb \\
 & $y=ax$ の $a$ を、\R{平方完成}{へい|ほう|かん|せい}で\R{求}{もと}められる。 & \pageref{bk:3-16:genten} & \bkcb & \bkcb \\
 & $y=ax+b$ を\R{共分散}{きょう|ぶん|さん}と\R{分散}{ぶん|さん}で\R{求}{もと}め、\R{残差}{ざん|さ}で\R{確}{たし}かめられる。 & \pageref{bk:3-16:ippan} & \bkcb & \bkcb \\
{\sffamily\bfseries\color{\bkTT}3-17} & \R{表}{ひょう}や\R{樹形図}{じゅ|けい|ず}で\R{数}{かぞ}えて、\R{確率}{かく|りつ}と\R{期待値}{き|たい|ち}を\R{求}{もと}められる。 & \pageref{bk:3-17:kazoeru} & \bkcb & \bkcb \\
 & パスカルの\R{三角形}{さん|かく|けい}で、\R{表}{おもて}が $k$ \R{枚}{まい}になる\R{数}{かず}を\R{求}{もと}められる。 & \pageref{bk:3-17:pascal} & \bkcb & \bkcb \\
 & \R{正規分布}{せい|き|ぶん|ぷ}の68\%と95\%を\R{使}{つか}える。 & \pageref{bk:3-17:seiki} & \bkcb & \bkcb \\
\end{longtblr}
\endgroup
